% !TeX root = ../../Book1.tex
\section{闭值域与可解性}
\label{b1:sec:C103}

对于有限维矩阵，值域总是闭的，所以一致线性方程组的极限仍然一致。
无穷维空间中的区别在于：右端可能收敛，所需原像却越来越大。
例如\chapref{b1:ch:08}的对角算子 \(Te_n=n^{-1}e_n\)，可以把单位向量压得任意小。
本节把排除这种现象的条件写成估计，并说明这个估计怎样传到对偶方程。

\subsection{去掉核以后的下界}
\label{b1:subsec:C10301}

\begin{theorem}[闭值域的估计判别]\label{b1:thm:C1-closed-range-estimate}
设 \(X,Y\) 为 Banach 空间，\(T\in\mathcal B(X,Y)\)。以下条件等价：
\begin{equivalences}
\item\label{b1:item:C1-range-closed}\(\operatorname{ran}T\) 在 \(Y\) 中闭；
\item\label{b1:item:C1-quotient-bound}存在 \(C>0\)，使
\[
 \operatorname{dist}(x,\ker T)\leq C\|Tx\|\quad(x\in X).
\]
\end{equivalences}
在这些条件下，对每个 \(y\in\operatorname{ran}T\)，方程 \(Tx=y\) 的解组成一个商类，
而从 \(\operatorname{ran}T\) 到 \(X/\ker T\) 的解映射有界。
\end{theorem}
\begin{proof}
若值域闭，则它是 Banach 空间。由%
\cref{b1:prop:C1-quotient-factor}，
\(\widetilde T:X/\ker T\rightarrow\operatorname{ran}T\) 是有界双射。
有界逆定理给出 \(\|\widetilde T^{-1}y\|\leq C\|y\|\)，即所述估计。

反过来，设估计成立，\(Tx_n\rightarrow y\)。对任意 \(m,n\)，有
\[
 \|Qx_n-Qx_m\|
 =\operatorname{dist}(x_n-x_m,\ker T)
 \leq C\|Tx_n-Tx_m\|.
\]
所以 \(Qx_n\) 在完备的商空间中收敛于某个 \(u\)。连续性给出
\(\widetilde Tu=y\)，因此 \(y\in\operatorname{ran}T\)。这里没有声称原来的
\(x_n\) 收敛；它们可以在核的方向任意漂移，真正收敛的是商类。
\end{proof}

当 \(\ker T=\{0\}\) 时，估计简化为 \(\|Tx\|\geq c\|x\|\)。
它比单射性强：单射性只禁止某个非零向量被压成零，估计则禁止单位向量列的像趋于零。
当核非零时，单位向量本来就可能在核内；正确的归一化是到核的距离等于 \(1\)，
而不是单纯要求 \(\|x\|=1\)。

闭值域保证有界的商类解映射，但不自动给出一个取值于 \(X\) 的有界线性选解算子。
要线性选取代表，需要核有一个闭补空间。有限维核总有这种补空间：取核的一组基，
将其坐标泛函用 Hahn--Banach 延拓到 \(X\)，所得有限秩投影的核就是所需补空间。
Hilbert 空间中则可以统一使用 \((\ker T)^\perp\)。

\subsection{Banach 闭值域定理}
\label{b1:subsec:C10302}

\begin{theorem}[Banach，闭值域]\label{b1:thm:C1-banach-closed-range}
设 \(X,Y\) 为 Banach 空间，\(T\in\mathcal B(X,Y)\)。
则 \(\operatorname{ran}T\) 范数闭，当且仅当 \(\operatorname{ran}T'\) 范数闭。
当它们闭时，
\[
 \operatorname{ran}T={}^{\circ}(\ker T'),\quad
 \operatorname{ran}T'=(\ker T)^\circ.
\]
\end{theorem}
\begin{proof}
先设 \(\operatorname{ran}T\) 闭，并取%
\cref{b1:thm:C1-closed-range-estimate}中的常数 \(C\)。
若 \(f\in(\ker T)^\circ\)，在值域上定义
\(\varphi(Tx)=f(x)\)。因 \(f\) 消去核，此定义与原像选择无关，而且
\[
 |\varphi(Tx)|=|f(x)|
 \leq\|f\|\operatorname{dist}(x,\ker T)
 \leq C\|f\|\cdot\|Tx\|.
\]
Hahn--Banach 定理把 \(\varphi\) 延拓为 \(g\in Y^*\)，于是 \(T'g=f\)。
反向包含 \(\operatorname{ran}T'\subset(\ker T)^\circ\) 不需闭性就成立。
因此 \(\operatorname{ran}T'=(\ker T)^\circ\)，后者是闭子空间。
值域的另一等式由\cref{b1:prop:C1-range-annihilator}去掉闭包得到。

现在设 \(\operatorname{ran}T'\) 闭。将闭值域的估计判别用于 \(T'\)，可取 \(C>0\) 使
\[
 \operatorname{dist}(g,\ker T')\leq C\|T'g\|\quad(g\in Y^*).
\]
记 \(Z=\overline{\operatorname{ran}T}\)。对 \(z\in Z\)，每个 \(h\in\ker T'\) 都消去 \(z\)，故
\begin{equation}\label{b1:eq:C1-dual-lifting-bound}
 |g(z)|\leq\|z\|\operatorname{dist}(g,\ker T')
 \leq C\|z\|\cdot\|T'g\|.
\end{equation}
这个不等式先给出范数可控的近似原像，再由级数把近似变为精确。

具体地，令 \(D=\overline{T(\overline B_X(0,1))}\)，这是 \(Y\) 中闭、凸且平衡的集合。
我们断言 \(\overline B_Z(0,1)\subset CD\)。
若某个 \(\|z\|\leq1\) 的 \(z\in Z\) 不属于 \(CD\)，则凸分离定理给出
\(g\in Y^*\)，使
\[
 \operatorname{Re}g(z)>\sup_{v\in CD}\operatorname{Re}g(v)=C\|T'g\|.
\]
最后的等式使用复平衡性；实标量情形使用正负对称性。
这与\eqref{b1:eq:C1-dual-lifting-bound}矛盾，断言得证。

取 \(z\in Z\)。由刚得到的包含关系，可以选 \(x_1\) 满足
\(\|x_1\|\leq C\|z\|\) 及 \(\|z-Tx_1\|\leq\|z\|/2\)；
若某次余项已经为零，后面全部取零。
否则对余项再次作同样选择，得到
\[
 \|x_j\|\leq C\,2^{-(j-1)}\|z\|,\quad
 \left\|z-T\sum_{j=1}^n x_j\right\|\leq2^{-n}\|z\|.
\]
由 \(X\) 完备，级数 \(x=\sum_jx_j\) 收敛，且 \(\|x\|\leq2C\|z\|\)。
连续性给出 \(Tx=z\)。因此 \(Z\subset\operatorname{ran}T\)，值域闭。
再用证明的前半部分，得到第二个值域等式。
\end{proof}

把\term*[banach-bizhiyu]{Banach 闭值域定理}[Banach closed range theorem]用于
存在性问题时，应分别核对“测试消去”与“值域闭”。
前者往往来自积分分部，后者则来自先验估计、紧性或某种有限维约化。
仅写出形式上的伴随方程，不足以证明原方程有解。

上述反向证明也解释了完备性的位置。对偶估计首先只保证任意小残差，
然后以几何级数控制校正项，原空间的完备性才允许把它们求和。
若直接从 \(z\in\overline{T(B_X)}\) 选出精确原像，就跳过了这一技术桥。

\subsection{Hilbert 空间中的正交解}
\label{b1:subsec:C10303}

\begin{corollary}\label{b1:cor:C1-hilbert-closed-range}
设 \(T:H\rightarrow G\) 为 Hilbert 空间间的有界线性算子。
若 \(\operatorname{ran}T\) 闭，则
\[
 \operatorname{ran}T=(\ker T^*)^\perp,\quad
 \operatorname{ran}T^*=(\ker T)^\perp.
\]
每个满足 \(y\perp\ker T^*\) 的右端有唯一的解 \(x_0\perp\ker T\)；
它是全部解中范数最小的一个，并满足 \(\|x_0\|\leq C\|y\|\)。
\end{corollary}
\begin{proof}
由 Riesz 对应 \(T'R_G=R_HT^*\)，Banach 闭值域定理中的消去条件变为正交条件。
任一解 \(x\) 可唯一分解为 \(x=x_0+n\)，其中 \(x_0\perp\ker T\)，\(n\in\ker T\)。
投影后仍有 \(Tx_0=y\)，两个这样的解相减同时属于核及其正交补，只能为零。
由勾股等式 \(\|x_0+n\|^2=\|x_0\|^2+\|n\|^2\)，最小范数性成立。
最后，\(\|x_0\|=\operatorname{dist}(x_0,\ker T)\)，所以前述估计适用。
\end{proof}

若右端 \(y\) 不满足正交条件，也可以先将其投影到 \(\operatorname{ran}T\)。
此时最小化 \(\|Tx-y\|\) 等价于解
\[
 T^*Tx=T^*y.
\]
这是因为 \(Tx-y\) 与值域正交当且仅当被 \(T^*\) 消去。
只要值域闭，投影存在，从而得到最小二乘解；若值域不闭，最小残差的下确界可能不取得。
闭值域因此同时联系精确可解性、解的稳定性和最小二乘问题。
